\chapter{Alcohol}

\section*{Definition and Overview}
Alcohol (ethanol) is a psychoactive substance produced by the fermentation of sugars, consumed worldwide in the form of beer, wine, spirits, and mixed drinks. It acts as a depressant on the central nervous system, meaning it slows brain activity. In small amounts alcohol produces relaxation and reduced inhibition; in larger amounts it impairs judgement, coordination, and consciousness.

Because alcohol is legal and widely available, its health effects are often underestimated. Understanding what alcohol is, how the body processes it, how drinking is measured, and what its short- and long-term effects are helps people make informed decisions about their intake.

\section*{Epidemiology}
Alcohol is the most widely used psychoactive substance in Australia and most other countries. A majority of adults drink alcohol at least occasionally, and a substantial minority drink at levels that increase their long-term health risk. Heavy drinking and binge drinking are most common in young adults.

Alcohol is a leading preventable cause of injury, illness, and premature death worldwide. It contributes to a large proportion of emergency department presentations, especially on weekends and public holidays, and is a major factor in road crashes, falls, and assaults.

\section*{Aetiology and Pathophysiology}
Alcohol is absorbed quickly from the stomach and small intestine into the bloodstream and is broken down mainly by the liver at a fairly constant rate. The effects depend on the blood alcohol concentration, which is influenced by how much is drunk, how quickly, body size, food in the stomach, and how tolerant the person has become.

\begin{itemize}
    \item \textbf{One standard drink:} contains about 10 grams of alcohol -- for example, a middy of full-strength beer, a small glass of wine, or a 30 ml measure of spirits
    \item \textbf{Effects on the brain:} alcohol enhances the action of the inhibitory neurotransmitter GABA and blocks glutamate, slowing brain activity
    \item \textbf{Low levels:} relaxation, reduced inhibition, and slightly slowed reaction time
    \item \textbf{Moderate levels:} slurred speech, poor coordination, and impaired judgement and memory
    \item \textbf{High levels:} vomiting, stupor, and eventually coma with depressed breathing (see the chapter on alcohol poisoning)
    \item \textbf{Metabolism:} the liver processes most of the alcohol consumed; the remainder is lost in breath, sweat, and urine
\end{itemize}

\section*{Clinical Features}
The immediate effects of alcohol depend on how much is consumed and how quickly.

\begin{itemize}
    \item Feeling relaxed and less inhibited
    \item Slowed reaction time, slurred speech, and unsteadiness
    \item Impaired judgement, leading to risk-taking behaviour
    \item Headache, nausea, and dehydration -- the hangover -- the next day
    \item Blackouts, in which parts of the drinking episode cannot be remembered
    \item At higher levels: vomiting, confusion, drowsiness, and unconsciousness
\end{itemize}

\section*{Differential Diagnosis}
When someone appears drunk or is behaving oddly, other causes must be considered.

\begin{itemize}
    \item Head injury or concussion, which can follow a fall while drinking
    \item Low blood sugar (hypoglycaemia), particularly in people with diabetes
    \item Drug overdose, such as from benzodiazepines or opioids, which can be mixed with alcohol
    \item Stroke, seizure, or meningitis, which can mimic intoxication
    \item Mental health conditions presenting with unusual behaviour
\end{itemize}

\section*{When to Seek Medical Care}
Most intoxication resolves with time, but severe intoxication can be dangerous and requires emergency care.

\textbf{Call 000 (or local emergency services) for:}
\begin{itemize}
    \item A person who is unconscious or cannot be roused
    \item Slow or irregular breathing (fewer than about eight breaths per minute)
    \item Vomiting while drowsy, or signs of choking
    \item Cold, clammy, or bluish skin, or a seizure
    \item A head injury, or a fall that may have injured the head
\end{itemize}

Also seek help if drinking is causing problems in your life, relationships, work, or health, or if you find it hard to control how much you drink.

\section*{Diagnosis and Investigations}
\begin{itemize}
    \item \textbf{History:} how much and how often the person drinks, and in what situations
    \item \textbf{Screening questionnaires:} such as the AUDIT, which identifies hazardous, harmful, or dependent drinking
    \item \textbf{Breath or blood alcohol testing:} measures the level of intoxication
    \item \textbf{Blood tests:} may include liver function tests, full blood count, and checks for iron, blood sugar, or cholesterol problems in heavy drinkers
    \item \textbf{Physical examination:} for signs of liver disease, nerve damage, or other effects of heavy drinking
\end{itemize}

\section*{Management}
\begin{itemize}
    \item \textbf{Low-risk drinking:} follow the national guidelines -- no more than ten standard drinks a week and no more than four on any one day
    \item \textbf{Alcohol-free days:} plan days each week without alcohol
    \item \textbf{Pacing:} drink slowly, with food, and alternate with water or non-alcoholic drinks
    \item \textbf{Avoiding harm:} never drink and drive, and avoid drinking when you need to concentrate or care for others
    \item \textbf{Cutting down:} measure drinks, choose lower-alcohol options, and set limits before going out
    \item \textbf{Seeking help:} see a GP if cutting down is difficult or if there are signs of dependence
\end{itemize}

\section*{Complications}
\begin{itemize}
    \item Injury, accidents, and falls while intoxicated
    \item Alcohol poisoning, which can be fatal
    \item Liver disease, pancreatitis, heart disease, and several cancers with long-term heavy use
    \item Dependence and withdrawal
    \item Harm to relationships, work, and finances
    \item Harm to others, including traffic accidents and violence
\end{itemize}

\section*{Prognosis and Outcomes}
The effects of alcohol depend on how much and how often a person drinks. Many of the short-term harms resolve when drinking stops, and reducing intake lowers the long-term risks. The health damage from long-term heavy drinking can be substantial, but some effects, such as early liver disease, may improve if drinking is stopped early enough. The outlook is best for people who recognise a problem early and reduce or stop their drinking.

\section*{Prevention and Self-Care}
\begin{itemize}
    \item \textbf{Know the guidelines:} understand the low-risk drinking limits and stay within them
    \item \textbf{Alcohol-free days:} include days each week without alcohol
    \item \textbf{Safe transport:} never drink and drive, and plan transport before you start drinking
    \item \textbf{Pacing:} drink with food, pace yourself, and have a glass of water between drinks
    \item \textbf{Pregnancy:} do not drink when pregnant or planning pregnancy, or when breastfeeding
    \item \textbf{Children:} keep alcohol out of reach of children, and do not supply it to people under 18
    \item \textbf{Medicines:} check medicine labels and avoid alcohol if the medicine warns against it
\end{itemize}

\section*{Sources and Further Reading}
The following sources provide reliable, up-to-date information on alcohol and its health effects.

\begin{itemize}
    \item NHS. \textit{Alcohol advice: effects, guidelines, and cutting down.} https://www.nhs.uk
    \item World Health Organization. \textit{Alcohol: global health information and policy.} https://www.who.int
    \item Centers for Disease Control and Prevention. \textit{Alcohol and public health.} https://www.cdc.gov
    \item Mayo Clinic. \textit{Alcohol use and health risks.} https://www.mayoclinic.org
    \item MedlinePlus. \textit{Alcohol -- health effects and safe use.} https://medlineplus.gov
\end{itemize}
